\documentclass[notheorems,hyperref={bookmarks=true},usepdftitle=false]{beamer}
\hypersetup{
pdfpagemode=FullScreen,
pdftitle={Xây dựng ứng dụng hỗ trợ tư vấn pháp luật sử dụng công nghệ trí tuệ nhân tạo},
pdfauthor={Vũ Văn Nghĩa - 20206205}
}
\usepackage[utf8]{vietnam}
\usepackage{mdframed}
\usepackage[english]{babel}
\usepackage{ulem}
\usepackage{pifont}
\usepackage[mathscr]{eucal}
\usepackage{amsfonts,amsmath, amsthm, amssymb,amsxtra,latexsym,amscd,graphics,graphpap}
\usepackage{rotating}
\usepackage{graphicx}
\usepackage{amsbsy}
\usepackage{epsfig}
\usepackage{natbib}
\usepackage{cases}

\usepackage{pdfpages}
\usepackage{multicol}
\usepackage{subfigure}
\usepackage{multimedia}
\usepackage{boxedminipage}
\usepackage{tikz}
\usepackage{nicefrac}
\setbeamertemplate{theorems}[numbered]
\usetheme{Boadilla}

\usecolortheme{rose}
\usefonttheme[onlylarge]{structurebold}
\usefonttheme{professionalfonts}
\setbeamercolor{title}{fg=red!80!black,bg=blue!20!white}
\setbeamertemplate{section in head/foot shaded}[default][20]
\setbeamertemplate{subsection in head/foot shaded}[default][20]
\setbeamertemplate{navigation symbols}{}
\setbeamertemplate{blocks}[rounded][shadow=true]
\setbeamerfont{small}{size=\small}
\setbeamerfont{footnote}{size=\footnotesize}
\setbeamerfont{script}{size=\scriptsize}
\setbeamerfont{tiny}{size=\tiny}

\theoremstyle{plain}
\renewcommand{\thetable}{\Roman{table}}

\mode<presentation>
\setbeamertemplate{navigation symbols}{}

\usepackage{wrapfig}
\usepackage{multicol}

\title[{\makebox[.15\paperwidth]{Khoa Toán - Tin}}]{Xây dựng ứng dụng hỗ trợ tư vấn pháp luật \\ sử dụng công nghệ trí tuệ nhân tạo}

\author[Vũ Văn Nghĩa - 20206205]{
Khoa Toán - Tin \\
Đồ án tốt nghiệp \\[0.2cm]
\begin{tabular}{ll}
Giảng viên hướng dẫn: & TS. Trần Anh Tú \\
Sinh viên thực hiện: & Vũ Văn Nghĩa \\
Mã số sinh viên: & 20206205 \\
Lớp: & Toán Tin 02 - K65
\end{tabular}
}
% \date{Hà Nội, \today}
% \date{Hà Nội, 07/2026}
% \date{\href{https://example.com}{Hà Nội, 07/2026}}
\date{\href{https://20206205tech.github.io/docs-slide}{Hà Nội, 07/2026}}




% học thuộc
% LLM... điểm yếu
% RAG

% Bổ sung
% mfa
% các sự kiện event

\renewcommand{\sfdefault}{cmss}
\renewcommand{\rmdefault}{cmr}
\renewcommand{\ttdefault}{cmtt}
\numberwithin{equation}{section}
\usepackage{array}
\usepackage{ragged2e}
\justifying
\usepackage{blindtext}
\renewcommand{\baselinestretch}{1.15}

\DeclareMathOperator*{\argmin}{arg\,min}
\newtheorem{prop}{Proposition}
\usepackage{amsmath}

\newtheorem{definition}{Definition}
\newtheorem{theorem}{Theorem}
\newtheorem{corollary}{Corollary}
\newtheorem{hypothesis}{Hypothesis}

\renewcommand{\thedefinition}{\arabic{section}.\arabic{definition}}
\providecommand{\thelemma}{}
\renewcommand{\thelemma}{\arabic{section}.\arabic{lemma}}
\renewcommand{\thetheorem}{\arabic{section}.\arabic{theorem}}

\usepackage{esvect}
\usepackage{booktabs}
\usepackage{multirow}

\usepackage{pgfpages}
\usepackage{makecell}

% \hypersetup{
% pdfpagelayout=TwoPageRight, % Chuyển sang trang bên phải khi nhấp chuột trái hoặc phải
% pdfpagemode=FullScreen, % Mở tài liệu ở chế độ toàn màn hình khi bắt đầu
% }
%%%%%%%%%%%%%%%%%%%%%%%%%%%%%%%%%%%%%%%%%%%

\newif\ifdev
\devtrue
% \devfalse

\ifdev

\setbeamercolor{background canvas}{bg=black} % Thiết lập nền toàn bộ slide thành màu đen
\setbeamercolor{normal text}{fg=white} % Thiết lập màu chữ mặc định là trắng

% \setbeameroption{show notes on second screen=right} % Bật chức năng hiển thị ghi chú ở màn hình bên phải (màn hình phụ)

\fi

\begin{document}
% \begin{footnotesize}

% Tiêu đề
\begin{frame}
\frametitle{}
\maketitle

\note{

Em xin chào thầy cô hội đồng

Em tên là Vũ Văn Nghĩa. MSSV 20206205

Hôm nay em xin trình bày về chủ đề đồ án tốt nghiệp

Xây dựng ứng dụng hỗ trợ tư vấn pháp luật
sử dụng công nghệ trí tuệ nhân tạo

được Hướng dẫn bởi TS Trần Anh Tú

Chuyển Slide

Nội dung báo cáo gồm có 100 phần \dots

Đọc các mục lên

}

\end{frame}

% \input{contents/Muc_luc_tu_dong}

% \section{Bài toán nghiệp vụ}

\input{contents/Su_can_thiet}

\input{contents/Danh_sach_Vai_tro}

\input{contents/So_do_Use_Case_Khach_vang_lai}
\input{contents/So_do_Use_Case_Nguoi_dung_co_ban}
\input{contents/So_do_Use_Case_Nguoi_dung_VIP}
\input{contents/So_do_Use_Case_Quan_tri_vien}

\section{Xây dựng hệ thống AI}
\subsection{Đường ống dữ liệu}

\input{contents/Duong_ong_du_lieu_van_ban_phap_luat_VBPL}
\input{contents/Quy_trinh_cac_buoc_thu_thap_du_lieu_van_ban_phap_luat}
\input{contents/Quy_trinh_cac_buoc_moi}
\input{contents/Duong_ong_du_lieu_Phap_dien}
\input{contents/Quy_trinh_cac_buoc_trong_duong_ong_du_lieu_Phap_dien}

\subsection{Agentic RAG với LangGraph}

\input{contents/Kien_truc_Agentic_RAG_bang_LangGraph}

\input{contents/Tong_quan_Kien_truc_Vi_dich_vu_trong_Du_an}

\section{Chi tiết các dịch vụ}

\input{contents/Kien_truc_Dich_vu_Thanh_toan_Payment_Service}

\input{contents/Kien_truc_Dich_vu_Tai_lieu_Document_Service}

\input{contents/Kien_truc_Dich_vu_Chatbot_va_Tro_chuyen}

\input{contents/Kien_truc_tich_hop_dam_thoai_thoi_gian_thuc_Voice_Agent}

\ifdev
\input{dev}
\fi

\section{Các công nghệ khác}

\input{contents/Quy_trinh_tu_dong_hoa_CICD_bang_GitHub_Actions}
\input{contents/Quan_ly_Ha_tang_duoi_dang_ma_IaC_bang_Terraform}
\input{contents/Cau_hinh_Kong_API_Gateway_bang_Declarative_Configuration}
\input{contents/docker}
\input{contents/Cau_hinh_tai_nguyen_Kubernetes_trong_GitOps_Repository}

\input{contents/Giao_dien_quan_ly_trien_khai_GitOps_bang_ArgoCD}

\section{Giao diện kết quả}

\input{contents/Giao_dien_Email_xac_thuc_dang_ky_tai_khoan}
\input{contents/Giao_dien_Thiet_lap_xac_thuc_hai_lop_MFA_qua_ma_QR}
\input{contents/Giao_dien_Tao_lien_ket_chia_se_cuoc_tro_chuyen}
\input{contents/Giao_dien_Dang_ky_goi_VIP_thanh_cong}
\input{contents/Giao_dien_Lua_chon_giong_noi_cua_tro_ly_ao}
\input{contents/Giao_dien_Dam_thoai_giong_noi_thoi_gian_thuc_voi_tro_ly_ao}

\section{Kết luận \& Hướng phát triển}
\input{contents/Ket_luan_cua_do_an}
\input{contents/Huong_phat_trien_trong_tuong_lai}

\section{Tài liệu tham khảo}

\input{contents/Tai_lieu_tham_khao}


\begin{frame}{Cảm ơn}
\centering
\Huge{Thanks for listening!}
% \Huge{Em cảm ơn thầy, cô đã lắng nghe!}
\end{frame}


\end{document}
